\section{Herramientas de búsqueda}
A continuación se presentan cinco herramientas destacadas para la búsqueda y/o organización de papers académicos en Internet:\\

\subsection{Semantic Scholar}
\textbf{Autor:} Allen Institute for AI (Ai2), fundado por Paul Allen.\\
\textbf{Año de publicación:} 2015.\\
\textbf{Enlace:} \url{https://www.semanticscholar.org/}.\\
\textbf{Descripción:} Motor de búsqueda AI para más de 200 millones de papers con resúmenes automáticos y recomendaciones. Se busca por keywords, autores o DOI; permite explorar grafos de citas, TLDR generados por IA y métricas como influencia; es ideal para descubrimiento rápido sin necesidad de login.
\subsection{Open Knowledge Maps}
\textbf{Autor:} Peter Kraker (fundador).\\
\textbf{Año de publicación:} 2016.\\
\textbf{Enlace:} \url{https://openknowledgemaps.org/}.\\
\textbf{Descripción:} Crea mapas visuales de conocimiento agrupando papers por temas clave vía procesamiento de lenguaje natural. Genera clústeres temáticos con papers representativos.
\subsection{Litmaps}
\textbf{Autor:} Axton Pitt (co-fundador y CEO).\\
\textbf{Año de publicación:} 2021.\\
\textbf{Enlace:} \url{https://app.litmaps.com/}.\\
\textbf{Descripción:} Permite visualizar mapas de literatura científica para encontrar papers relevantes por citas, fecha y conectividad.
\subsection{Connected Papers}
\textbf{Autor:} Connected Paper Inc.\\
\textbf{Año de publicación:} 2020.\\
\textbf{Enlace:} \url{https://www.connectedpapers.com/}.\\
\textbf{Descripción:} Sirve para visualizar grafos de similitud entre papers científicos basados en citas y contenido. Se usa ingresando un DOI, título o URL de un paper inicial; genera un mapa interactivo donde nodos cercanos son similares, permitiendo explorar trabajos relacionados, filtrar por año o citas, y exportar listas.
\subsection{ResearchRabbit}
\textbf{Autor:} ResearchRabbit\\
\textbf{Año de publicación:} 2021.\\
\textbf{Enlace:} \url{https://www.researchrabbit.ai/}.\\
\textbf{Descripción:} Ayuda a descubrir literatura mediante mapas visuales de conexiones entre papers, autores y temas. Se ingresa un paper semilla o palabras clave, permite agregar notas, sincronizar con Zotero y recibir alertas de nuevos papers relevantes.